%% ---- 4. Método --------------------------------------------------------------
\section{Método}

\subsection{Estimação por Monte Carlo}
\label{sec:mc}

\paragraph{Distribuição de fundo.}
Cada posição de um fundo é um item sorteado independentemente com
probabilidade proporcional à sua popularidade no treino. Nunca sorteamos $v$,
seus 50 vizinhos mais próximos no espaço de \emph{embeddings} (similaridade do
cosseno) nem os itens que o próprio usuário consumiu, para que um
preenchimento ``neutro'' não devolva evidência a favor de $v$. O fundo
representa um histórico qualquer, sem relação com $v$, e essa escolha define o
que ``suficiente'' significa; por isso a testamos com duas ablações
(Seção~\ref{sec:ablacao}).

\paragraph{Estimadores.}
Para cada instância $(s,v)$ sorteamos uma única vez $M$ fundos, de
$r^{(1)}$ a $r^{(M)}$, e estimamos
\begin{equation}
  \hat N(z) = \frac{1}{M} \sum_{m=1}^{M} \delta\big(v \mid s[X \leftarrow r^{(m)}]\big)
  \label{eq:Nhat}
\end{equation}
\begin{equation}
  \hat S(z) = -\,\frac{1}{M} \sum_{m=1}^{M} \delta\big(v \mid s[\bar X \leftarrow r^{(m)}]\big)
  \label{eq:Shat}
\end{equation}
$\widehat{\PN}$ e $\widehat{\PS}$ são as frações dos mesmos $M$ sorteios com
$\delta<0$ e $\delta\ge0$, respectivamente. Cada parcela é uma passada do
SASRec, de modo que avaliar uma máscara custa $2M$ passadas.

\paragraph{Fundos fixos.}
O mesmo conjunto de fundos serve a todas as máscaras e a todos os métodos da
instância (números aleatórios comuns). A aptidão vira uma função
determinística de $z$: o NSGA-II não compara indivíduos contaminados por
ruídos diferentes, e máscaras repetidas podem ser respondidas de um
\emph{cache} sem custo. O preço é um viés de seleção: a busca pode se ajustar
aos $M$ fundos que viu. Por isso as frentes finais são reavaliadas com
$M'=100$ fundos independentes, e todos os resultados reportados usam essa
reavaliação.

\subsection{NSGA-II sobre máscaras}

\begin{figure}[t]
\centering
\begin{tikzpicture}[font=\scriptsize, >=stealth, node distance=5mm and 4mm,
  box/.style={draw, rounded corners=1.5pt, align=center, inner sep=2.5pt,
              minimum height=8mm}]
  \node[box] (pop) {População\\[-1pt]$\{z_k\}$};
  \node[box, right=of pop] (ver) {Verificador: substitui\\[-1pt]itens por $M$ fundos e\\[-1pt]consulta o SASRec};
  \node[box, right=of ver] (obj) {$(\hat N,\hat S,C)$\\[-1pt]de cada $z_k$};
  \node[box, below=of obj] (sel) {Ordenação não\\[-1pt]dominada e\\[-1pt]aglomeração};
  \node[box, below=of ver] (var) {Seleção, cruzamento\\[-1pt]e mutação};
  \draw[->] (pop) -- (ver);
  \draw[->] (ver) -- (obj);
  \draw[->] (obj) -- (sel);
  \draw[->] (sel) -- (var);
  \draw[->] (var) -| node[pos=0.25, below, font=\tiny] {nova geração} (pop);
  \node[font=\scriptsize, below=2mm of var, align=center]
        {Ao final: máscaras não dominadas entre todas as avaliadas};
\end{tikzpicture}
\caption{Visão geral. Cada máscara $z_k$ gera $2M$ históricos com
substituições, avaliados em lote pelo SASRec.}
\label{fig:visao}
\end{figure}

O indivíduo é a máscara $z$ (Figura~\ref{fig:visao}). Usamos o
NSGA-II~\cite{deb2002nsga2} da pymoo~\cite{blank2020pymoo} com:
\begin{itemize}
  \item \textbf{inicialização esparsa}: cada bit ligado com probabilidade
        $5/n$, para começar perto das explicações pequenas;
  \item \textbf{cruzamento uniforme} e \textbf{mutação por inversão de bit} com
        taxa $1/n$;
  \item \textbf{reparo}: uma máscara vazia recebe um bit aleatório;
  \item eliminação de indivíduos duplicados.
\end{itemize}
A frente devolvida é o conjunto não dominado de \emph{todas} as máscaras que a
execução avaliou, e não apenas a população final. O tamanho da população e $M$
foram ajustados em usuários separados dos de avaliação
(Seção~\ref{sec:ajuste}), e cada escolha de operador foi testada isoladamente
(Seção~\ref{sec:operadores}).

\paragraph{Variante semeada.}
O ranking por ablação individual (Seção~\ref{sec:base}) converge muito rápido
no início da busca. A variante \textbf{NSGA-II semeado} gasta primeiro o que
esse ranking gasta (cada posição sozinha e depois as $k$ melhores posições por
$\hat N$ e por $\hat S$, no máximo $3n$ máscaras). A população inicial são as
máscaras, entre essas, que a própria sobrevivência do NSGA-II manteria
(ordenação não dominada e aglomeração); o restante do orçamento é o NSGA-II
usual. Essa variante foi proposta depois da comparação principal e é
reportada à parte (Seção~\ref{sec:semeado}).

\subsection{Linhas de base}
\label{sec:base}

Todos os métodos consultam o mesmo verificador, com os mesmos fundos, e sua
frente é o conjunto não dominado de tudo o que avaliaram.
\begin{itemize}
  \item \textbf{AG escalarizado}: algoritmo genético de objetivo único com os
        mesmos operadores do NSGA-II. A aptidão é a soma ponderada dos três
        objetivos, cada um dividido pelo seu valor na máscara cheia ($N$, $S$)
        ou por $n$ ($C$). O orçamento é dividido igualmente entre 10 vetores de
        pesos (reticulado do simplex), uma execução por vetor, com população
        20. É o controle que separa o efeito de ser evolutivo do efeito de ser
        multiobjetivo.
  \item \textbf{Guloso}: sobre $\hat N + \hat S$ normalizados, inclui uma
        posição por vez a partir da máscara vazia e, com o que sobra do
        orçamento, remove uma por vez a partir da cheia.
  \item \textbf{Ranking por ablação}: avalia cada posição sozinha e depois os
        conjuntos das $k$ melhores posições por $\hat N$ e por $\hat S$, para
        todo $k$.
  \item \textbf{Recência} e \textbf{similaridade}: as $k$ posições mais
        recentes, ou os $k$ itens mais semelhantes a $v$, para todo $k$.
  \item \textbf{Aleatório}: subconjuntos de tamanho uniforme em $1..n$.
  \item \textbf{Enumeração}: todas as máscaras com uma ou duas posições, sem
        limite de orçamento, como referência exata para $|X|\le2$.
  \item \textbf{GECE}~\cite{scarcelli2025gece}, pela implementação do
        TRACE~\cite{kim2026trace}, sem alterações de código e com seus
        parâmetros padrão (população 100, taxa de mutação $0{,}1$). A condição
        de validade passada a ele é a nossa: $v$ sair do top-10 dos candidatos,
        com no máximo 10 posições editadas. Como o GECE edita a sequência,
        mapeamos cada sequência que ele consulta para o conjunto de posições
        em que ela difere do histórico, e avaliamos esses conjuntos com o
        nosso verificador. Essa implementação para no primeiro contrafactual
        válido; a variante \textbf{GECE com reinícios} a reinicia até gastar o
        orçamento inteiro.
\end{itemize}
